\scanpage{160}
\mspage{76}
% archive sheet number: 159

% slanted
\margin{\unsure{?} \uline{Démonstration \unsure{concluante|conclut}} \unsure{en vérité|en réalité} que si $\pi^{\pi_0}=\pi$ !!}

Prouvons c), i.e. que pour tt $\pi'>\pi\geq\pi_0$ et \deleted{tout} $x\in\operatorname{Ob}\operatorname{Filt}^{\pi'}$ on peut écrire

\[
\begin{aligned}
x=\varinjlim_I x_i\qquad & I\ \text{grand devant }\pi\\
&\text{les }x_i\text{ dans }\operatorname{Filt}^{\pi}\\
&\operatorname{Card} I\leq\pi'^{\pi}.
\end{aligned}\tag*{(4.3.6.5)}
\]

La démonstration se fait comme dans \deleted{Pour ceci, on peut supposer que} loc. cit. (p. \unsure{151|159}--152), en changeant légèrement la définition de $I$, comme désignant l'ens des sous-catégories pleines \inserted{$i$} de $S_{/x}$ qui sont filtrantes, de cardinal $\leq\pi$ \uline{et de plus grandes devant $\pi_0$}. (On pourra supposer \inserted{les} $S_{/x}$ filtrantes, en \deleted{\unsure{ajoutant}} \unsure{augmentant}\deleted{\unsure{t}} $S$ de façon à ce que $S$ soit stable \inserted{dans $M$} par \deleted{petites} $\varinjlim$ finies, donc les catégories $S_{/x}$ aussi, ce qui implique qu'elles sont filtrantes. Il est

% slanted
\margin{\unsure{?} Il n'est pas sûr que $I$ soit grand devant $\pi$, sauf si on suppose $\pi^{\pi_0}=\pi$ !}

On a donc prouvé $\mathrm{PTA}\Longrightarrow\mathrm{c})$,

\uline{Prouvons} $\mathrm c\Longrightarrow\mathrm d$. Soit \deleted{\unsure{Soit} $c\geq\pi_0$, et} \deleted{$c\geq\pi_0$, et soit}

\[
\pi_1=2^c\quad\text{avec }c\geq\pi_0.
\]

On sait alors que

\[
\begin{aligned}
\operatorname{Filt}^{(\pi_1)}M=M_{\pi_1}\overset{\mathrm{def}}{=}\ &\text{sous-catégorie \inserted{pleine de $M$} des}\\
&\text{objets }\pi_1\text{-accessibles}\\
&\text{de }M
\end{aligned}
\]

(loc. cit 9.16), ce qui implique notamment que $M_{\pi_1}$ est stable par $\varinjlim$ de cardinal $\leq\pi_1$, \inserted{en vertu de loc cit. 9.9. (p\deleted{\unsure{ag}} 145).}

\deleted{En effet, si $x=\varinjlim_{i\in I}x_i$, $\operatorname{card}\operatorname{Fl}I\leq\pi$, et $x_i\in M_{\pi_1}$, alors le foncteur \inserted{$h_x:M\unsure{\text{---}|\longrightarrow}\mathrm{Ens}$} représenté \unsure{par} \unsure{$h$|$x$} est la $\varprojlim_I$ des foncteurs $h_{x_i}$ ; comme ils sont $\pi_1$-accessibles, il en est de $=$ de $h_x$ (par loc. cit. 9.\unsure{6} (ii\unsure{i})).}

De plus, on

